% verdicts and plain-words explanations of the slide deck (hand-written; every number in them is a macro generated by
% tests/t18_numbers.py from the measured JSONs, or is quoted from a log named next to it)
\newcommand{\PlanarFound}{\yes}
\newcommand{\MethaneFound}{\yes}
\newcommand{\RaoSixtyOneFound}{\yes}
\newcommand{\AxialFound}{\yes}
\newcommand{\BellFound}{\yes}
\newcommand{\BellIdealFound}{\partly}
\newcommand{\MigdalFound}{\yes}
\newcommand{\PLANARVERDICT}{Four starts --- two distant ones, and two with the no-fold rule on (one starting outside it)
--- land on one value (\PlanarPairJ\ppm\ apart) and one wall (\PlanarPairWall\,mm per m), in 5--13 steps, 2--4\,min.
At the finest grid the result is \PlanarBFine\ppm\ from the bound, the closest 12-point curve \PlanarFitFine\ppm; in the
limit the result sits \PlanarBMinusFitLim\ppm\ from that curve. (Old driver: 27--28 segments with a radius floor 16
times finer than its default --- with the default it stopped 1\ppm\ short; with the rule on it stalled at $-49$\ppm, and
from a start breaking the rule it did not move.)}
\newcommand{\PLANARWHY}{The exact wall bends abruptly where the first lip wave lands. A smooth curve through 12
points cannot follow that bend: its wall angle is off by $\delta\theta\approx2\times10^{-3}$ rad on average, and a wall
angle error costs about $2\,\delta\theta^2\approx9$\ppm\ of thrust. No optimiser can beat the curve it is given; the
walls of the result and of the exact plug differ by $\sim$1\,mm per m where the thrust does not care (flat top).}
\newcommand{\METHANEVERDICT}{\yes. Two starts land on one point (\MethanePairJ\ppm,
walls \MethanePairWall\,mm per m apart) in \MethaneATrials\ and \MethaneCTrials\ steps: \MethaneAMinusFitWalk\ppm\ from
the closest curve of Rao's contour on the optimiser's grid, \MethaneAMinusFitLim\ppm\ in the limit. Same landing as the old
driver; this is the RDE-like plug ($\theta_\mathrm{wall}+\mu\ge5.5^\circ$, far from the steep regime of Rao 1961).}
\newcommand{\METHANEWHY}{Rao's method gives the best plug for exactly this inflow,
length and end point; the optimiser, which knows none of that, finds the same wall to a fraction of a millimetre per
metre. Above Rao on the coarse grid, below it on finer grids: the grid effect of every case, not a gain.}
\newcommand{\RAOVERDICT}{\yes, with the default settings. Both starts land on one point
(\RaoSixtyOnePairJ\ppm, walls \RaoSixtyOnePairWall\,mm per m apart) in \RaoSixtyOneBumpTrials\ and
\RaoSixtyOneHalfTrials\ steps: \RaoSixtyOneBumpMinusFitWalk\ppm\ from the closest curve of Rao's contour on the optimiser's grid,
\RaoSixtyOneBumpMinusFitLim\ppm\ in the limit; the wall is 0.4\,mm per m from Rao's. With the old driver the same starts
stopped \RaoSixtyOneOldBumpMinusFitWalk\ and \RaoSixtyOneOldHalfMinusFitWalk\ppm\ from that curve; it got there only with
hand-set short steps and two more hours per walk.}
\newcommand{\RAOWHY}{Rao's $\varepsilon=3.81$ plug is so steep near the throat that the
characteristics at the wall point slightly downward, and the thrust varies along a narrow ridge. The old driver walked
across the ridge and kept falling off it; the new one measures the ridge's shape once, keeps it, and walks along it.}
\newcommand{\AXIALVERDICT}{\yes. Two starts land on one point (\AxialPairJ\ppm,
walls \AxialPairWall\,mm per m apart) in \AxialBumpTrials\ and \AxialHalfTrials\ steps: \AxialBumpMinusFitWalk\ppm\ from
the closest curve of Rao's contour on the optimiser's grid, \AxialBumpMinusFitLim\ppm\ in the limit.}
\newcommand{\AXIALWHY}{The flow leaves this plug uniform and axial at ambient pressure: the most thrust any plug
can give. Rao's contour delivers it (to the consistency of GENO's data) and the optimiser finds Rao's contour; the
small remaining gap is the 9-point curve's and the grid's.}
\newcommand{\BELLVERDICT}{\yes. Two starts on opposite sides of the optimum (attachment \BellThBStartTic\textdegree\ and
\BellThBStartSteep\textdegree, thrust $-2.2\times10^{-3}$ and $-3.1\times10^{-2}$) end on one wall and one value. Optimised
on the grid twice as fine: attachment \BellThBTicEightOne\textdegree\ and \BellThBSteepEightOne\textdegree\ (closest curve
to GENO's \BellThBFit\textdegree), walls \BellPairWallFine\,mm per m apart, \BellPairJFine\ppm\ apart; in the limit
\BellTicEightOneMinusFitLim\ppm\ from the closest curve. Optimised on the coarser grid the same starts end
\BellTicMinusFitWalk\ppm\ above that curve on their own grid but \BellTicMinusFitLim\ppm\ in the limit: that grid is too
coarse to rank walls at the ppm level.}
\newcommand{\BELLWHY}{Rao's bell is the optimum; ours --- a curve through 8 points --- follows GENO's within
0.6\,mm per m. The optimiser sees the thrust only through its grid, and that grid's small error changes with the wall by a
few ppm: it gets used. One grid halving divides that effect by about 10 (from $-3.4$ to $+0.3$\ppm\ here).}
\newcommand{\BELLIDEALVERDICT}{\partly. From Rao's contour for length 6 followed by a cylinder the optimiser finds the
ideal bell: on the grid twice as fine \BellIdealExtendedEightOneMinusFitWalk\ppm\ from its closest curve on its own grid,
\BellIdealExtendedEightOneMinusFitLim\ppm\ in the limit. From the same contour stretched to the full length (attachment
9.1\textdegree) it converges --- correctly, its curvature measured there --- to a \textbf{second top}: attachment
\BellIdealThBStretchedEightOne\textdegree\ instead of \BellIdealThBFit\textdegree, wall 3\,\% of the throat radius away,
\BellIdealStretchedEightOneMinusFitLim\ppm\ in the limit. The two-start check flags it (\BellIdealPairJFine\ppm\ apart).}
\newcommand{\BELLIDEALWHY}{At its full length the ideal bell's thrust is extremely flat: walls up to 3\,\% of the throat
radius apart differ by only about 45\ppm. Both results are strict local maxima (all curvatures negative, every probing march lower) and between
them the thrust dips by 14\ppm\ (backup slide): two real tops of the computed thrust; a local optimiser stops on the one
it climbs. Not the 8-point wall (with 12 points the same start
still does not get there), not the throat-arc panel the objective leaves out, not the axial-Mach margin; the cause is
open. Two distant starts are the protection, and here they flag it.}
\newcommand{\MIGDALVERDICT}{\yes. With and without the no-fold rule --- the walk with the rule starting \emph{outside} it --- the two walks land on one point (\MigdalPairJ\ppm, walls \MigdalPairWall\,mm per m apart) in \MigdalFreeTrials\ and
\MigdalClsTrials\ steps. On the optimiser's grid \MigdalFreeWalk\ppm\ from GENO's perfect nozzle, on the finest grid
\MigdalFreeFine\ppm\ (limit \MigdalFreeMinusFitLim\ppm). The production walker (Newton metric, no-fold rule on, its own radius
rule) lands on the same point, $C_F$ equal to 9 digits, in 3 steps. Old driver: from a start breaking the rule it did not move.}
\newcommand{\MIGDALWHY}{The perfect nozzle turns all the flow axial at the exit; walls that do the same within the
grid's error give the same thrust. The result and GENO's walls differ by 0.5\,mm per m on the plug and 0.1 on the shroud,
where the thrust cannot tell them apart; the ppm left are the grid's, as in every case.}
\newcommand{\OPENONE}{\textbf{The grid sets the precision.} The walk maximises the thrust its grid computes and
uses that grid's small error: 1--2\ppm\ on the plug grids, $\sim$3--5\ppm\ on the bell engine at $41\times120$, more on the
flat ideal bell. One grid halving cuts the effect by about 10 (bell TOC: $-3.4\to+0.3$\ppm). Rule: walk on a grid whose
error varies with the wall less than the precision wanted, and re-march the result once on a twice finer grid. When the
grid's thrust is not smooth at the scale of the last steps the driver now says so and stops (``RESOLUTION''); that
stop was added when the coarse ideal-bell walk climbed grid steps, and replayed on the logs of every other walk here it
never fires.}
\newcommand{\OPENTWO}{\textbf{Two tops.} A local optimiser stops on the top it climbs: the ideal bell at its full
length has a second top 30--45\ppm\ lower (a 14\ppm\ valley between the two results); its cause is not identified (not the
wall description, not the throat panel, not the axial-Mach margin). Protection: always two distant starts; a third when
they disagree. Every other case here has one top (all starts agree to $<0.3$\ppm).}
\newcommand{\OPENTHREE}{\textbf{Starts and cost.} The bell engine refuses a conical start (internal shock,
outside the model) and long shallow starts (a cell's axial Mach below 1 near the throat). The walker's generic start on the
annular nozzle folds before the driver is called (the walker's own start logic, unchanged); the driver itself now restores
starts that break the no-fold rule. \textbf{Cost:} on the Rao posings one gradient takes 40--100\,s (today's s2 was loaded:
our runs got 25--35\,\% of a core), and the two measured curvatures (2 $\times$ 18 gradients) are $\sim$90\,\% of a walk:
1--1.7\,h today. One-sided differences would halve that, a compiled replay would cut it further; not done, so that the
runs shown here match the code that is now the default. Other cases: 2--45\,min per walk.}
\newcommand{\BELLVSAH}{$+0.3$\ppm\ (optimised on the grid twice as fine) to $-3.5$\ppm\ (coarser grid)}
\newcommand{\HthFound}{\ifdefined\HthDoneFlag\yes\else\yes$^*$\fi}
\newcommand{\HTHVERDICTFINAL}{\yes. From a characteristic of Rao's own field, Humphreys' method gives a contour ``almost
identical to Rao's'' ($-2$\,lbf); our optimiser, from two starts, lands \HthBumpFine\ and \HthHalfFine\ppm\ from Rao's
contour on the finest grid (\HthBumpMinusFitLim\ppm\ from its closest curve in the limit), walls \HthPairWall\,mm per m
apart.}
\newcommand{\HTHVERDICTPROV}{\yes\ on the optimiser's grid. Their Table 3 lies on GENO's Rao contour within
$5\times10^{-5}R_E$ on the designed part; Humphreys' method from a characteristic of Rao's field: $-2$\,lbf
($-58$\ppm). Our SQP on the same problem: \HthProv. The check on finer grids is still running.}
\newcommand{\HTHVERDICT}{\ifdefined\HthDoneFlag\HTHVERDICTFINAL\else\HTHVERDICTPROV\fi}
\newcommand{\HTHWHY}{Same inflow, length and end point: one best plug. Rao builds it from the exit characteristic,
Humphreys from the given start line, our optimiser by climbing on thrust; all three find it. Their design 4\,\% lower
is for their own engine (another start line, wall friction): a different problem.}
